\documentclass[12pt]{standalone}
\usepackage{fontspec}
\setsansfont{TeX Gyre Heros}
\usepackage{tikz}
\usetikzlibrary{arrows.meta}
% =====================================================================
% 统一符号层（2026-10-20）
% 用法：在 fig_<号>.tex 的 preamble 里 % =====================================================================
% 统一符号层（2026-10-20）
% 用法：在 fig_<号>.tex 的 preamble 里 % =====================================================================
% 统一符号层（2026-10-20）
% 用法：在 fig_<号>.tex 的 preamble 里 \input{../_symbols/symbols.tex}
% 依赖：fontspec（已 \setsansfont{TeX Gyre Heros}）、graphicx 不需要
% 规则见同目录《符号验收清单.md》
% =====================================================================

% ---- 字体角色（全部走 sans = TeX Gyre Heros）----
\newcommand{\SIt}[1]{\textsf{\textit{#1}}}          % 斜体变量
\newcommand{\SBf}[1]{\textsf{\textbf{#1}}}          % 粗体
\newcommand{\SBfIt}[1]{\textsf{\textbf{\textit{#1}}}} % 粗斜体变量
\newcommand{\SNum}[1]{\textsf{\textup{#1}}}         % 直立数字

% ---- 上下标：字号/位置写死，避免数学样式与手调混用 ----
% 审计结论（2026-10-20）：源图下标/底字高比在 0.77–0.92em，0.72em 偏小。
% 默认 0.80em；仍需按源图逐图标定（44.4/44.5 约 0.85–0.92，39/41/46 约 0.77–0.80）。
\newcommand{\FSup}[2]{#1\raisebox{0.30em}{\fontsize{0.72em}{0.82em}\selectfont #2}}
\newcommand{\FSub}[2]{#1\raisebox{-0.30em}{\fontsize{0.80em}{0.92em}\selectfont #2}}

% ---- 分数：数字直立、方头粗杆；杆宽 = max(分子宽,分母宽) + 2pt ----
% 杆厚 4pt（配合 20–26pt 标签）；小字号图可局部调 0.17em
% 2026-10-20 修正：原 \shortstack 负 kern 不随字号缩放（23pt 下总高仅 30pt、杆贴字），
% 改为 offinterlineskip + em 定间距：上间隙 0.12em、下间隙 0.13em
% （37.2 源图 23pt 实测四标签上/下间隙均值 2.6/4.0pt；取近平衡且总高保持 42–43pt）。
\newcommand{\FFrac}[2]{%
  \begingroup
  \setbox0=\hbox{#1}\setbox2=\hbox{#2}%
  \dimen0=\wd0 \ifdim\wd2>\dimen0 \dimen0=\wd2 \fi
  \advance\dimen0 by 2pt
  \leavevmode\vbox{\offinterlineskip
    \hbox to \dimen0{\hss\box0\hss}%
    \kern0.12em
    \hbox{\rule{\dimen0}{4pt}}%
    \kern0.13em
    \hbox to \dimen0{\hss\box2\hss}%
  }%
  \endgroup}

% ---- 真减号（审计：源图减号 15–18px 宽、3–5px 厚；Heros 常规字重偏细）----
% 可选参数 = 宽度（em）：正文 0.58；上标内用 \FMinus[0.72]
\newcommand{\FMinus}[1][0.58]{\rule[0.30em]{#1em}{0.13em}}

% ---- 直撇（审计：44.2 用弯引号、44.3 用直撇，必须统一为直楔撇）----
\newcommand{\FPrime}{\textsf{\textquotesingle}}

% ---- 乘号（审计：46.2 的 × 比源小 25–30%，标签内需按源校准）----
% 用 \FTimes 时在 node 内与底字同字号，并在目检时量宽高。


% ---- 斜杠分数（审计 A：35.1 的 r/3 用 38pt 斜杠手搓）----
\newcommand{\FSlashFrac}[2]{#1\raisebox{-0.18em}{\fontsize{1.35em}{1.35em}\selectfont\textsf{/}}\kern-0.06em #2}

% ---- 交集号 ∩（Heros 无此字，固定回退字号/基线，避免逐图不同）----
\newcommand{\FIntersect}{\raisebox{0.02em}{\fontsize{0.95em}{1em}\selectfont\mfallback{∩}}}

% ---- 括号/开区间（图形层，不是字形）----
% 一律 \draw ... arc[start angle=, end angle=, radius=]：
%   · 端点记号背离它所界的区间；
%   · 左端/右端分别 10× 特写核对；
%   · 闭区间压在轴上时竖脊右移半个线宽；
%   · 同一张图同类记号同一种写法。

% 依赖：fontspec（已 \setsansfont{TeX Gyre Heros}）、graphicx 不需要
% 规则见同目录《符号验收清单.md》
% =====================================================================

% ---- 字体角色（全部走 sans = TeX Gyre Heros）----
\newcommand{\SIt}[1]{\textsf{\textit{#1}}}          % 斜体变量
\newcommand{\SBf}[1]{\textsf{\textbf{#1}}}          % 粗体
\newcommand{\SBfIt}[1]{\textsf{\textbf{\textit{#1}}}} % 粗斜体变量
\newcommand{\SNum}[1]{\textsf{\textup{#1}}}         % 直立数字

% ---- 上下标：字号/位置写死，避免数学样式与手调混用 ----
% 审计结论（2026-10-20）：源图下标/底字高比在 0.77–0.92em，0.72em 偏小。
% 默认 0.80em；仍需按源图逐图标定（44.4/44.5 约 0.85–0.92，39/41/46 约 0.77–0.80）。
\newcommand{\FSup}[2]{#1\raisebox{0.30em}{\fontsize{0.72em}{0.82em}\selectfont #2}}
\newcommand{\FSub}[2]{#1\raisebox{-0.30em}{\fontsize{0.80em}{0.92em}\selectfont #2}}

% ---- 分数：数字直立、方头粗杆；杆宽 = max(分子宽,分母宽) + 2pt ----
% 杆厚 4pt（配合 20–26pt 标签）；小字号图可局部调 0.17em
% 2026-10-20 修正：原 \shortstack 负 kern 不随字号缩放（23pt 下总高仅 30pt、杆贴字），
% 改为 offinterlineskip + em 定间距：上间隙 0.12em、下间隙 0.13em
% （37.2 源图 23pt 实测四标签上/下间隙均值 2.6/4.0pt；取近平衡且总高保持 42–43pt）。
\newcommand{\FFrac}[2]{%
  \begingroup
  \setbox0=\hbox{#1}\setbox2=\hbox{#2}%
  \dimen0=\wd0 \ifdim\wd2>\dimen0 \dimen0=\wd2 \fi
  \advance\dimen0 by 2pt
  \leavevmode\vbox{\offinterlineskip
    \hbox to \dimen0{\hss\box0\hss}%
    \kern0.12em
    \hbox{\rule{\dimen0}{4pt}}%
    \kern0.13em
    \hbox to \dimen0{\hss\box2\hss}%
  }%
  \endgroup}

% ---- 真减号（审计：源图减号 15–18px 宽、3–5px 厚；Heros 常规字重偏细）----
% 可选参数 = 宽度（em）：正文 0.58；上标内用 \FMinus[0.72]
\newcommand{\FMinus}[1][0.58]{\rule[0.30em]{#1em}{0.13em}}

% ---- 直撇（审计：44.2 用弯引号、44.3 用直撇，必须统一为直楔撇）----
\newcommand{\FPrime}{\textsf{\textquotesingle}}

% ---- 乘号（审计：46.2 的 × 比源小 25–30%，标签内需按源校准）----
% 用 \FTimes 时在 node 内与底字同字号，并在目检时量宽高。


% ---- 斜杠分数（审计 A：35.1 的 r/3 用 38pt 斜杠手搓）----
\newcommand{\FSlashFrac}[2]{#1\raisebox{-0.18em}{\fontsize{1.35em}{1.35em}\selectfont\textsf{/}}\kern-0.06em #2}

% ---- 交集号 ∩（Heros 无此字，固定回退字号/基线，避免逐图不同）----
\newcommand{\FIntersect}{\raisebox{0.02em}{\fontsize{0.95em}{1em}\selectfont\mfallback{∩}}}

% ---- 括号/开区间（图形层，不是字形）----
% 一律 \draw ... arc[start angle=, end angle=, radius=]：
%   · 端点记号背离它所界的区间；
%   · 左端/右端分别 10× 特写核对；
%   · 闭区间压在轴上时竖脊右移半个线宽；
%   · 同一张图同类记号同一种写法。

% 依赖：fontspec（已 \setsansfont{TeX Gyre Heros}）、graphicx 不需要
% 规则见同目录《符号验收清单.md》
% =====================================================================

% ---- 字体角色（全部走 sans = TeX Gyre Heros）----
\newcommand{\SIt}[1]{\textsf{\textit{#1}}}          % 斜体变量
\newcommand{\SBf}[1]{\textsf{\textbf{#1}}}          % 粗体
\newcommand{\SBfIt}[1]{\textsf{\textbf{\textit{#1}}}} % 粗斜体变量
\newcommand{\SNum}[1]{\textsf{\textup{#1}}}         % 直立数字

% ---- 上下标：字号/位置写死，避免数学样式与手调混用 ----
% 审计结论（2026-10-20）：源图下标/底字高比在 0.77–0.92em，0.72em 偏小。
% 默认 0.80em；仍需按源图逐图标定（44.4/44.5 约 0.85–0.92，39/41/46 约 0.77–0.80）。
\newcommand{\FSup}[2]{#1\raisebox{0.30em}{\fontsize{0.72em}{0.82em}\selectfont #2}}
\newcommand{\FSub}[2]{#1\raisebox{-0.30em}{\fontsize{0.80em}{0.92em}\selectfont #2}}

% ---- 分数：数字直立、方头粗杆；杆宽 = max(分子宽,分母宽) + 2pt ----
% 杆厚 4pt（配合 20–26pt 标签）；小字号图可局部调 0.17em
% 2026-10-20 修正：原 \shortstack 负 kern 不随字号缩放（23pt 下总高仅 30pt、杆贴字），
% 改为 offinterlineskip + em 定间距：上间隙 0.12em、下间隙 0.13em
% （37.2 源图 23pt 实测四标签上/下间隙均值 2.6/4.0pt；取近平衡且总高保持 42–43pt）。
\newcommand{\FFrac}[2]{%
  \begingroup
  \setbox0=\hbox{#1}\setbox2=\hbox{#2}%
  \dimen0=\wd0 \ifdim\wd2>\dimen0 \dimen0=\wd2 \fi
  \advance\dimen0 by 2pt
  \leavevmode\vbox{\offinterlineskip
    \hbox to \dimen0{\hss\box0\hss}%
    \kern0.12em
    \hbox{\rule{\dimen0}{4pt}}%
    \kern0.13em
    \hbox to \dimen0{\hss\box2\hss}%
  }%
  \endgroup}

% ---- 真减号（审计：源图减号 15–18px 宽、3–5px 厚；Heros 常规字重偏细）----
% 可选参数 = 宽度（em）：正文 0.58；上标内用 \FMinus[0.72]
\newcommand{\FMinus}[1][0.58]{\rule[0.30em]{#1em}{0.13em}}

% ---- 直撇（审计：44.2 用弯引号、44.3 用直撇，必须统一为直楔撇）----
\newcommand{\FPrime}{\textsf{\textquotesingle}}

% ---- 乘号（审计：46.2 的 × 比源小 25–30%，标签内需按源校准）----
% 用 \FTimes 时在 node 内与底字同字号，并在目检时量宽高。


% ---- 斜杠分数（审计 A：35.1 的 r/3 用 38pt 斜杠手搓）----
\newcommand{\FSlashFrac}[2]{#1\raisebox{-0.18em}{\fontsize{1.35em}{1.35em}\selectfont\textsf{/}}\kern-0.06em #2}

% ---- 交集号 ∩（Heros 无此字，固定回退字号/基线，避免逐图不同）----
\newcommand{\FIntersect}{\raisebox{0.02em}{\fontsize{0.95em}{1em}\selectfont\mfallback{∩}}}

% ---- 括号/开区间（图形层，不是字形）----
% 一律 \draw ... arc[start angle=, end angle=, radius=]：
%   · 端点记号背离它所界的区间；
%   · 左端/右端分别 10× 特写核对；
%   · 闭区间压在轴上时竖脊右移半个线宽；
%   · 同一张图同类记号同一种写法。

% 乘号：源图 C×z0 与 f(C×z0) 两处均 16×16，基底字约 31pt ⇒ 42pt 同写法
\newcommand{\FTimes}{\raisebox{-2pt}{\fontsize{42}{46}\selectfont\textsf{\textbf{×}}}}
\begin{document}
\begin{tikzpicture}[x=1pt, y=-1pt, line cap=round, line join=round]

  % ===== labels: one label = one node, drawn first, white background =====
  \node[anchor=base west, inner sep=0pt, fill=white, font=\fontsize{30.5}{36}\selectfont\bfseries] at (40,79) {\SBfIt Z};
  \node[anchor=base west, inner sep=0pt, fill=white, font=\fontsize{29}{34}\selectfont\bfseries] at (14,196) {\SBfIt W};
  \node[anchor=base west, inner sep=0pt, fill=white, font=\fontsize{30.7}{37}\selectfont\bfseries] at (88.5,189.5) {\FSub{\SBfIt z}{\SNum 0}};
  \node[anchor=center, inner sep=0pt, fill=white, font=\fontsize{31}{37}\selectfont\bfseries] at (291,424) {\SBfIt C};
  \node[anchor=center, inner sep=0pt, fill=white, font=\fontsize{30.5}{36}\selectfont\bfseries] at (483,425.5) {\SBfIt X};
  \node[anchor=base west, inner sep=0pt, fill=white, font=\fontsize{30.7}{37}\selectfont\bfseries] at (505.5,175) {\SBfIt C\FTimes\FSub{\SBfIt z}{\SNum 0}};
  \node[anchor=base west, inner sep=0pt, fill=white, font=\fontsize{30.5}{36}\selectfont\bfseries] at (510.5,311) {\FSup{\SBfIt f}{\FMinus\SNum 1}\SBf{(}\SBfIt U\SBf{)}};
  \node[anchor=center, inner sep=0pt, fill=white, font=\fontsize{30}{36}\selectfont\bfseries] at (660,45) {\SBfIt f};
  \node[anchor=center, inner sep=0pt, fill=white, font=\fontsize{30.5}{36}\selectfont\bfseries] at (937.5,44) {\SBfIt U};
  \node[anchor=center, inner sep=0pt, fill=white, font=\fontsize{29}{34}\selectfont\bfseries] at (1049.5,29) {\SBfIt Y};
  \node[anchor=base west, inner sep=0pt, fill=white, font=\fontsize{31}{37}\selectfont\bfseries] at (955.5,231) {\SBfIt f\SBf{(}\SBfIt C\FTimes\FSub{\SBfIt z}{\SNum 0}\SBf{)}};

  % ===== hatch: upper band, "/" family (x+y=const) =====
  \begin{scope}
    \clip (216,55) rectangle (363.5,145);
    \foreach \x in {49.2,67.5,...,342.0} { \draw[line width=3pt] (\x,145) -- (\x+90,55); }
  \end{scope}

  % ===== hatch: lower band, "/" family =====
  \begin{scope}
    \clip (216,232) rectangle (363.5,396);
    \foreach \x in {50.2,68.4,...,359.6} { \draw[line width=3pt] (\x,396) -- (\x+164,232); }
  \end{scope}

  % ===== cross hatch band ("/" and "\") first; blob covers it =====
  \begin{scope}
    \clip (216,145) rectangle (363.5,232);
    \foreach \x in {23.0,41.3,...,352.4} { \draw[line width=3pt] (\x,232) -- (\x+87,145); }
    \foreach \x in {109.8,128.1,...,366.0} { \draw[line width=3pt] (\x,145) -- (\x+87,232); }
  \end{scope}

  % ===== blob f^-1(U): one closed outline on top of all hatch（源图不填白）=====
  \draw[line width=5pt] plot[smooth cycle, tension=0.5] coordinates {
    (216,136) (235,127) (250,114) (266,107) (282,106) (296,113) (307,123) (318,132)
    (331,133) (345,127) (357,121) (367,117) (385,118) (405,126) (422,136) (436,146) (445,158)
    (452,172) (457,186) (458,202) (456,221) (451,233) (444,243) (435,251) (425,257)
    (413,261) (402,261) (390,259) (378,256) (367,253) (359,251) (340,255) (321,258)
    (307,264) (295,270) (284,284) (276,290) (266,296) (254,297) (240,298) (216,296)
    (210,296) (197,290) (185,282) (174,272) (168,259) (167,245) (167,232) (168,218)
    (167,204) (163,190) (157,177) (152,164) (153,150) (160,136) (172,133) (185,132) (200,131)};

  % ===== band boundaries =====
  \draw[line width=5pt, line cap=butt] (216,145) -- (363.5,145);
  \draw[line width=5pt, line cap=butt] (216,232) -- (363.5,232);

  % ===== C x z0 thick segment =====
  \draw[line width=7pt, line cap=butt] (216,182) -- (363.5,182);

  % ===== strip boundaries and outer square Z =====
  \draw[line width=5pt] (216,55) -- (216,396);
  \draw[line width=5pt] (363.5,55) -- (363.5,396);
  \draw[line width=5pt] (72,55) rectangle (494,396);

  % ===== thick C bar on bottom edge =====
  \draw[line width=9pt, line cap=butt] (216,401) -- (364,401);

  % ===== W brace（与 C×z0 的 y 区间 145–232 对齐）=====
  \draw[line width=3pt] (61,145) .. controls (51,147) and (53,154) .. (53,163)
    -- (53,180) .. controls (53,186) and (49,189) .. (43,191)
    .. controls (49,193) and (53,196) .. (53,202)
    -- (53,220) .. controls (53,229) and (51,233) .. (61,232);
  % 区间端点短横：与大括号两端同 y，跨过矩形边（x=72）
  \draw[line width=3.5pt, line cap=butt] (60,145) -- (84,145);
  \draw[line width=3.5pt, line cap=butt] (60,232) -- (84,232);

  % ===== z0 dot =====
  \fill (72,181.5) circle (7.5);

  % ===== kidney Y outer =====
  \draw[line width=5pt] plot[smooth cycle, tension=0.5] coordinates {
    (990,402) (956,398) (927,382) (905,356) (892,326) (887,293) (882,259) (865,231)
    (840,207) (810,194) (778,185) (758,160) (756,126) (767,95) (786,67) (813,47)
    (842,32) (873,22) (906,16) (939,15) (972,21) (1003,32) (1032,46) (1060,65)
    (1086,86) (1108,112) (1123,141) (1130,173) (1132,208) (1130,242) (1126,275)
    (1118,307) (1102,336) (1079,361) (1052,381) (1022,394)};

  % ===== inner U =====
  \draw[line width=4.5pt] plot[smooth cycle, tension=0.5] coordinates {
    (872,192) (853,190) (838,180) (829,165) (826,147) (828,129) (836,113) (848,99)
    (862,86) (877,76) (894,70) (911,65) (929,62) (948,62) (966,66) (984,70)
    (1000,78) (1016,86) (1029,100) (1036,117) (1040,135) (1038,154) (1032,170)
    (1018,182) (1000,186) (983,181) (967,171) (955,161) (944,157) (932,160)
    (921,169) (910,178) (899,185) (887,189)};

  % ===== arc f(C x z0) with end dots =====
  \draw[line width=6pt] (861,146) .. controls (895,104) and (940,103) .. (975,110);
  \fill (861,146) circle (8.5);
  \fill (975,110) circle (8.5);

  % ===== arrows =====
  \draw[-{Triangle[length=20pt,width=14pt]}, line width=4pt] (568,95) .. controls (610,66) and (690,63) .. (748,93);
  \draw[-{Triangle[length=20pt,width=13pt]}, line width=4pt] (509,163) -- (368,173);
  \draw[-{Triangle[length=17pt,width=12pt]}, line width=4pt] (503,298) .. controls (487,292) and (466,278) .. (440,256);
  \draw[-{Triangle[length=21pt,width=14.5pt]}, line width=4pt] (953,210) .. controls (936,194) and (922,182) .. (919,160)
    .. controls (917,144) and (918,126) .. (919,112);

  % ===== pin canvas =====
  \pgfresetboundingbox
  \path[use as bounding box] (0,0) rectangle (1149,449);
\end{tikzpicture}
\end{document}
